\begin{frame}{Caractéristiques principales du NoSQL}

    \begin{itemize}
        \item Dénormalisation des données
        \item Scalabilité et partitionnement
        \item Haute disponibilité
        \item Cohérence des données
        \item Indexation
        \item Transactions
        \item Schéma flexible
    \end{itemize}
\end{frame}

\begin{frame}{Dénormalisation}

	\begin{definitionbox}
		La \textbf{dénormalisation} est un processus de structuration des données dans une base de données NoSQL pour améliorer les performances de lecture et de requêtage, en introduisant une certaine redondance dans les données.
	\end{definitionbox}

	La dénormalisation est le processus inverse de la normalisation. Elle consiste à regrouper les données dans une même table ou document, même si cela entraîne une redondance, afin de réduire le nombre de jointures nécessaires pour accéder aux données.

\end{frame}
